\section{User requests}
\label{sec:requests}

A Request asks the Mint to evaluate a name change. A Respond returns an OTP for a pending Request. Neither is a Name Note and neither binds a transition to a note commitment.

\subsection{Request encoding}

User-to-Mint memos are a 512-byte field. A parser removes the maximal trailing sequence of zero bytes. The remaining bytes MUST be UTF-8. The colon byte (\texttt{0x3a}) separates fields. The first field MUST be the exact ASCII bytes \texttt{ZNS}. Extra fields, missing fields, and empty required fields are rejected. The \texttt{<years?>} slot of an \action{update} is the one optional field: it may be empty.

The name field MUST contain between 1 and 63 bytes, each of which is ASCII \texttt{a}--\texttt{z} or \texttt{0}--\texttt{9}.

A Request has one of the following forms:

\begin{center}
\texttt{ZNS:claim:<term>:<name>:<ua>}\\
\texttt{ZNS:update:<years?>:<name>:<ua>}\\
\texttt{ZNS:release:<name>:<ua>}
\end{center}

The \texttt{<term>} field of a \action{claim} is either the exact ASCII bytes \texttt{forever} or \texttt{<N>y}, where \texttt{<N>} is a canonical decimal integer from 1 to 99 with no leading zeroes. \texttt{forever} requests a registration without fixed expiration. The \texttt{<years?>} field of an \action{update} is either empty, which carries the current \field{expires\_at} forward unchanged, or \texttt{<N>y} under the same spelling rules, which requests an extension of \texttt{<N>} years. A \action{claim} MUST NOT use \texttt{none}, and an \action{update} MUST NOT use \texttt{none}: the slot is left empty or carries years.

A Respond has one of the following forms:

\begin{center}
\texttt{ZNS:update:<name>:<ua>:<otp>}\\
\texttt{ZNS:release:<name>:<ua>:<otp>}
\end{center}

A Request or Respond MUST include a non-empty Unified Address. A \action{claim} is only a Request. Request and Respond memos do not encode \field{expires\_at} or \field{prev\_rcm}.

An \action{update} or \action{release} Request is the first step of the authorization procedure specified in Section~\ref{sec:authorization}. A Respond is the third step: the controller sending the OTP back to the Mint. The \texttt{<otp>} field is exactly six ASCII decimal digits, including leading zeroes.

A parser for these memos MUST reject any memo that does not match one of the forms above, including a Relay memo as specified in Section~\ref{sec:authorization}.
